\section{Introduction}
\label{sec:intro}

Acoustic Side-Channel Attacks (ASCAs) exploit the acoustic emanations from electronic devices to infer sensitive information without directly accessing the targeted systems. This adversarial method represents a substantial cybersecurity risk, particularly today as microphones become ubiquitous components integrated into common consumer devices, such as smartphones, laptops, digital assistants, and conferencing equipment. Although initially perceived as a novel theoretical threat, practical acoustic attacks have emerged to compromise sensitive information in various contexts ~\cite{zhuang2009keyboard, carrara2015acoustic, compagno2017don, mehrnezhad2018stealing}.

Previous research has demonstrated the feasibility of recovering typed passwords and PINs covertly by analyzing keystroke sounds~\cite{zhuang2009keyboard, asonov2004keyboard, berger2006dictionary, anand2018keyboard, harrison2023practical}, reconstructing confidentially printed text from dot-matrix printer noises in sensitive environments such as medical and banking facilities~\cite{backes2010acoustic}, and extracting cryptographic information or CPU operation details from subtle hardware acoustic signals~\cite{shamir2004acoustic, ninan2024second}. Among these, keyboard-focused acoustic attacks have been proven to be especially threatening due to their widespread applicability and ability to compromise passwords and confidential communications effectively~\cite{zhuang2009keyboard, asonov2004keyboard, harrison2023practical, halevi2012closer}.

Early ASCA methodologies primarily depended on straightforward statistical analyses and signal-processing techniques, leveraging acoustic features like cepstrum coefficients and inter-keystroke timing patterns for keystroke recovery~\cite{zhuang2009keyboard}. Despite enhancements through cross-correlation analysis and time-based localization techniques that improved keystroke identification~\cite{berger2006dictionary, de2019differential}, these approaches exhibited significant limitations when faced with realistic scenarios involving ambient noise. Machine-Learning (ML) methods such as Hidden Markov Models (HMMs) and Support Vector Machines (SVMs) subsequently improved accuracy by modeling sequential structures and key feature distinctions~\cite{backes2010acoustic,wang2016accurate}.

However, these early systems remained vulnerable to real environmental conditions, where significant amounts of ambient noise and recording artifacts could drastically degrade results—potentially causing incorrect keystroke identifications. Previous research has shown accuracy reductions of 30–50\% in noisy environments~\cite{halevi2012closer, harrison2023practical}. Our experiments confirm this, with accuracy dropping over 50\% under high ambient noise. Such inaccuracies are critical, particularly when attackers aim to recover sensitive information such as passwords, where even minor incorrect inferences make full password reconstruction near-impossible. Recent Deep-Learning (DL) approaches, notably Convolutional Neural Networks (CNNs) and architectures like CoAtNet~\cite{dai2021coatnet}, have further advanced acoustic classification accuracy by enabling effective hierarchies of acoustic feature extraction~\cite{harrison2023practical}. 

Despite these latest improvements, the fundamental challenge--vulnerability to noisy environmental conditions--remains largely unaddressed in the existing literature. Moreover, incorrect keystroke predictions under such conditions render traditional ASCA pipelines significantly less practical and reliable. Therefore, solutions for the noisy problem are still warranted. Our key observation and insight about it is that solving this problem requires either: \textbf{(i)} an increased model's capacity of inferring contextual information from longer sequences, allowing the model to learn that an initially noisily typed word is the same as a futurely collected non-noisy word, making the model learn a noisy representation for the characters directly from noisy data; or \textbf{(ii)} an approach to fix misidentified information from the contexts, exploiting the fact that one does not type random words, but the ones that best fit the conversation context. In this case, we would train the model with non-noisy data and allow the model to make a wrong prediction from the noisy spectrogram, a spectrogram ``typo'', but we would fix the prediction based on the context. However, technical solutions to deploy these strategies were not available in the literature until recently.

At the same time we identified the above limitations and needs, transformer architectures have emerged as a powerful class of neural network solutions capable of capturing complex long-range correlations and contextual relationships within data. Originally developed for sequence modeling in the language domain, transformers have demonstrated remarkable performance across multiple modalities—including language (LLMs), vision (VTs), and more recently audio—by leveraging their ability to globally correlate data patterns. In particular, vision transformers (e.g., ViT~\cite{dosovitskiy2020image} and Swin~\cite{liu2021swin}) successfully model spatial correlations in images and spectrogram representations, while large language models (e.g., GPT-3~\cite{brown2020language} and GPT-4~\cite{achiam2023gpt}) excel at contextual inference, error detection, and text correction tasks. These advantages naturally position transformers as candidate solutions for deploying our above-mentioned strategies. Despite these advantages, to the best of our knowledge, no research has yet leveraged transformer architectures—either vision or language variants—to handle noisy acoustic keystroke data. This position our work as the first-of-its-kind in making ASCAs practical.

Considering the above, we postulate three key scientific hypothesis: \textbf{(1)} The global modeling capability of visual transformers, particularly their unparalleled skill in capturing long-range contextual relationships, might effectively mitigate errors introduced by acoustic noise and improve ASCA reliability, making them the most suitable candidate for solving the problem via the first above-pointed strategy \textit{(i)}; and, similarly \textbf{(2)} Large Language Models (LLMs) have significant error correction (e.g., rewritting, passphrase, translation, and so on) capabilities, which makes them the most suitable solutions to solve the problem via the second above-mentioned strategy \textit{(ii)}; finally, \textbf{(3)} These solutions can complement each other and be used in tandem to solve the problem at a higher scale.

To test our hypothesis in practice, we developed a novel transformer-driven framework for enhancing acoustic side-channel attacks on keyboards and conducted experiments with it on reference datasets. More specifically, we leveraged advanced transformer-based architectures, employing both VTs (for spectrogram image classification) and LLMs (for contextual error correction), to substantially improve accuracy and robustness in noisy real-world settings. Our framework will be made available as open-source. \footnote{\url{https://github.com/Botacin-s-Lab/EchoCrypt}}

Our experimental evaluation was conducted using the Phone (keystrokes recorded via a smartphone microphone) and Zoom (keystrokes captured through Zoom audio call) datasets ~\cite{harrison2023practical}, considering two realistic settings. To simulate noise conditions, we introduced low, medium, and high noise levels. We compared GPT-4o~\cite{hurst2024gpt} with Llama-3.2-3B (fine-tuned) and tested Llama-3.2-1B, 3B, and 8B (non-fine-tuned)~\cite{touvron2023llama} to assess their effectiveness in mitigating noise-related errors. 

In sum, the contributions in this work are the following:

\begin{itemize}
    \item We design, implement, and evaluate VT-based frameworks for ASACs based on the Swin transformer and the CoAtNet model and demonstrate how their hypothesized capabilities of identifying long-range patterns enable them to establish the new SOTA performance in the reference ASAC benchmark.
    \item We design, implement, and evaluate LLM-based frameworks for ASACs based on GPT-4 and demonstrate that their hypothesized correction abilities allow extending the previous VT operation to noisy scenarios.
    \item We demonstrate that fine-tuned lightweight LLMs (via Low-Rank Adaptation, LoRA) can achieve typo-correction performance comparable to very large LLMs like GPT-4o, while having substantially fewer parameters—making these solutions practical even for resource-constrained attack scenarios.
\end{itemize}

In sum, the main findings of this work are the following:

\begin{enumerate}
   \item Our proposed VT framework achieved a new SOTA performance on the reference dataset, with an increase in absolute precision of 5.0\% on the Phone dataset and 5.9\% on the Zoom dataset compared to the previous CNN baselines.
   \item Our proposed LLM framework, incorporating GPT-4o to handle noisy cases, when applied to the \textit{EnglishTense} dataset (1K sentences), significantly boosts BLEU scores from approximately 0.07 to 0.90 under intermediate noise conditions, transforming an impractical classification case into a case with almost perfect key recovery.
   \item Our proposed fine-tuned LLM (Llama-3.2-3B), via Low-Rank Adaptation (LoRA), achieves accuracy scores closely comparable to GPT-4o (98-99\% of it) while being 67x smaller, which makes ASACs not only practical in noisy scenarios but also brings ASACs to one's pocket.
\end{enumerate}

The remainder of this paper is structured as follows. Section~\ref{sec:motiv} presents a clear example of our solution's goals and potential. Section~\ref{sec:background} presents background information to support our developments towards making ASCAs practical. Section \ref{sec:method} details our methodology, including our proposed vision transformer and LLM approaches, data preprocessing strategies, noise simulation methods, fine-tuning procedures, and experimental setups. Section \ref{sec:evaluation} presents our evaluation results, highlighting comparative analysis between previous state-of-the-art methods and our approach across two datasets and multiple metrics, and demonstrating the effectiveness of our proposed transformer architectures under noisy conditions. Section \ref{sec:discussion} provides an in-depth discussion of our results, practical implications, and limitations, outlining directions for future research. Section~\ref{sec:related} positions our contributions among related work, highlighting previous methodologies in acoustic keyboard side-channel classification, the evolution of classification techniques, and the emergence of transformer-based methods. Finally, Section \ref{sec:conclusion} concludes the paper by summarizing our key contributions and findings.